%\subsection{1  Лекция 1: Дифференциальные уравнения в частных производных второго порядка}



Многие задачи механики сплошной среды и физики приводят к исследованию дифференциальных уравнений с частными производными второго порядка (ДУЧП2П). Так, например.



1) При изучении различных видов волн (упругих, звуковых, электромагнитных) мы приходим к \textit{волновому уравнению}:
\[
	\frac{\partial ^{2} u}{\partial t^{2} } =c^{2} \left(\frac{\partial ^{2} u}{\partial x^{2} } + \frac{\partial ^{2} u}{\partial y^{2} } + \frac{\partial ^{2} u}{\partial z^{2} } \right) + f(x,y,z,t)
\] 
где $c$ -- скорость распространения волн в данной среде.

\bigskip

2) Процессы распространения тепла в однородном изотропном теле описываются \textit{уравнением теплопроводности}
\[
\frac{\partial u}{\partial t} =a^{2} \left(\frac{\partial ^{2} u}{\partial x^{2} } +\frac{\partial ^{2} u}{\partial y^{2} } +\frac{\partial ^{2} u}{\partial z^{2} } \right)+f(x,y,z,t)\] 
где $a^{2} $ -- коэффициент температуропроводности данной среды.

\bigskip

3) При рассмотрении установившегося теплового состояния в однородном изотропном теле мы приходим к \textit{уравнению Пуассона}:

\[\frac{\partial ^{2} u}{\partial x^{2} } +\frac{\partial ^{2} u}{\partial y^{2} } +\frac{\partial ^{2} u}{\partial z^{2} } =-\tilde{f}(x,y,z)\]
где   $\tilde{f}(x,y,z)=\frac{f(x,y,z)}{a^{2} } $

\bigskip

\noindent При отсутствии источников тепла внутри тела уравнение теплопроводности переходит в \textit{уравнение Лапласа}:
\[\frac{\partial ^{2} u}{\partial x^{2} } +\frac{\partial ^{2} u}{\partial y^{2} } +\frac{\partial ^{2} u}{\partial z^{2} } =0\] 

\bigskip

Волновое уравнение (ВУ), уравнение теплопроводности (УТ), уравнение Пуассона (УП) и уравнение Лапласа (УЛ) часто называют \textit{основными уравнениями математической физики}.

\bigskip

Каждое из основных уравнений математической физики имеет бесчисленное множество частных решений. При решении конкретной физической задачи, необходимо из всех этих решений выбрать то, которое удовлетворяет некоторым дополнительным условиям, вытекающим из её физического смысла. Итак, \textit{задачи математической физики состоят в отыскании решений уравнений в частных производных, удовлетворяющим некоторым дополнительным условиям}. Такими дополнительными условиями чаще всего являются так называемые \textit{граничные }(\textit{краевые})\textit{ условия}, т.е. условия, заданные на границе рассматриваемой среды, и \textit{начальные условия}, относящегося к одному какому-нибудь моменту времени, с которого начинается изучение данного физического явления.

\bigskip

Обратим внимание на одно важное обстоятельство.

\textit{Задача математической физики считается поставленной корректно }(\textit{в смысле Адамара})\textit{, если решение задачи, удовлетворяющее всем её условиям, существует, единственно и устойчиво} (непрерывно зависит от данных (параметров) задачи, т.е. малые изменения любого из данных задачи должно вызывать соответственно малое изменение решения). Требование устойчивости выражает физическую определенность поставленной задачи.

\bigskip

Дадим необходимые определения.

\begin{definition}
Пусть $G$ --- область в $\mathit{R}^{n} $, $x=(x_{1} ,...,x_{n} )$ -- точка из $G$. Дифференциальным уравнением в частных производных второго порядка (ДУЧП2П) с независимыми переменными $x_{1} $,{\dots},$x_{n} $ называется соотношение, связывающее значение неизвестной функции $u(x)=u(x_{1} ,...,x_{n} )$ и её частных производных первого и второго порядков: 
$\frac{\partial u}{\partial x_{1} } $, {\dots}, $\frac{\partial u}{\partial x_{n} } $, $\frac{\partial ^{2} u}{\partial x_{1}^{2} } $, {\dots},$\frac{\partial ^{2} u}{\partial x_{i} \partial x_{j} } $, {\dots}, $\frac{\partial ^{2} u}{\partial x_{n}^{2} } $
вида: 
\[F\left(x_{1} ,...,x_{n} ,u,\frac{\partial u}{\partial x_{1} } ,...,\frac{\partial u}{\partial x_{n} } ,\frac{\partial ^{2} u}{\partial x_{1}^{2} } ,...,\frac{\partial ^{2} u}{\partial x_{i} \partial x_{j} } ,...,\frac{\partial ^{2} u}{\partial x_{n}^{2} } \right)=0\] 
\end{definition}


\begin{definition}
Решением ДУЧП2П называется функция $u(x)$ класса $C^{2} (G)$, которая удовлетворяет данному уравнению.  
\end{definition}



\begin{definition}
Уравнение называется \textit{линейным} относительно старших производных -- производных второго порядка (ДУЧП2ПЛОСП), если оно имеет вид:
\[\sum _{i,j=1}^{n}a_{ij} (x)\frac{\partial ^{2} u}{\partial x_{i} \partial x_{j} }  +\Phi \left(x,u,{\rm grad}u\right)=0\] 
где коэффициенты при старших (вторых) производных $a_{ij} (x)\in C(G)$ ($i,j=1,...,n$) являются функциями только независимых переменных $x_{1} $,{\dots},$x_{n} $; а ${\rm grad}u=\left(\frac{\partial u}{\partial x_{1} } ,...,\frac{\partial u}{\partial x_{n} } \right)$ -- градиент функции $u(x)=u(x_{1} ,...,x_{n} )$.
\end{definition}


\noindent 
\begin{definition}
Коэффициенты при старших (вторых) производных образуют квадратную матрицу $A(x)=(a_{ij} (x))$ -- \textit{матрицу старших коэффициентов}:
\[A(x)=\left(\begin{array}{cccc} {a_{11} (x)} & {a_{12} (x)} & {...} & {a_{1n} (x)} \\ {a_{21} (x)} & {a_{22} (x)} & {...} & {a_{2n} (x)} \\ {.} & {.} & {...} & {.} \\ {a_{n1} (x)} & {a_{n1} (x)} & {...} & {a_{nn} (x)} \end{array}\right)\] 
\end{definition}

Для дальнейшего нам будет удобно считать матрицу $A(x)$ симметрической ($A^{T} (x)=A(x)$, $a_{ij} (x)=a_{ji} (x)$). Это не умаляет общности, так как в силу справедливых для функций из $C^{2} (G)$ равенств
\[\frac{\partial ^{2} u}{\partial x_{i} \partial x_{j} } =\frac{\partial ^{2} u}{\partial x_{j} \partial x_{i} } ,  i,j=1,...,n,\] 
\[\sum _{i,j=1}^{n}a_{ij} (x)\frac{\partial ^{2} u}{\partial x_{i} \partial x_{j} }  =\sum _{i,j=1}^{n}\frac{a_{ij} (x)+a_{ji} (x)}{2} \frac{\partial ^{2} u}{\partial x_{i} \partial x_{j} }  \] 
и, следовательно, не меняя уравнения, в качестве матрицы $A(x)$ можно взять симметрическую матрицу $\frac{1}{2} (A(x)+A^{T} (x))$ -- симметрическую часть матрицы $A(x)$:
\[A(x)=\frac{1}{2} (A(x)+A^{T} (x))+\frac{1}{2} (A(x)-A^{T} (x))\] 


\noindent 

\begin{note}
Если функция $\Phi \left(x,u,{\rm grad}u\right)$ линейна относительно аргументов $u$, $\frac{\partial u}{\partial x_{1} } $, {\dots}, $\frac{\partial u}{\partial x_{n} } $, то уравнение (ДУЧП2ПЛОСП) называется (просто) \textit{линейным} (без указания относительно чего).
\end{note}

\noindent 

Рассмотрим в области $G\subset \mathit{R}^{n} $ \textit{линейное дифференциальное уравнение в частных производных второго порядка} (ЛДУЧП2П):
\[Lu=\sum _{i,j=1}^{n}a_{ij} (x)\frac{\partial ^{2} u}{\partial x_{i} \partial x_{j} }  +\sum _{i=1}^{n}b_{i} (x)\frac{\partial u}{\partial x_{i} }  +c(x)u=f(x)\] 
где $a_{ij} (x)$, $b_{i} (x)$, $c(x)$ ($i,j=1,...,n$) -- заданные функции из $C(G)$ -- \textit{коэффициенты} уравнения, $f(x)$ -- заданная функция из $C(G)$ -- \textit{свободный член} (\textit{правая часть}) уравнения. Если $f(x)=f(x_{1} ,...,x_{n} )\equiv 0$, то уравнение (ЛДУЧП2П) называется \textit{однородным}, в противном случае -- \textit{неоднородным}. 



Здесь $L$ -- линейный дифференциальный оператор второго порядка:
\[L=\sum _{i,j=1}^{n}a_{ij} (x)\frac{\partial ^{2} }{\partial x_{i} \partial x_{j} }  +\sum _{i=1}^{n}b_{i} (x)\frac{\partial }{\partial x_{i} }  +c(x),\] 
преобразующий функции из $C^{2} (G)$ в функции из $C(G)$.



Если же все коэффициенты $a_{ij} $, $b_{i} $, $c$ ($i,j=1,...,n$) постоянны, то уравнение (ЛДУЧП2П) называется \textit{линейным уравнением с постоянными коэффициентами }(ЛДУЧП2ППК).

\bigskip 

Рассмотрим в области $G\subset \mathit{R}^{n} $ \textit{линейное дифференциальное уравнение в частных производных второго порядка с постоянными коэффициентами}\textbf{ }(ЛДУЧП2ППК):
\[Lu=\sum _{i,j=1}^{n}a_{ij} \frac{\partial ^{2} u}{\partial x_{i} \partial x_{j} }  +\sum _{i=1}^{n}b_{i} \frac{\partial u}{\partial x_{i} }  +cu=f(x)\] 
где $a_{ij} $, $b_{i} $, $c$ ($i,j=1,...,n$) -- \textit{коэффициенты} уравнения, $f(x)$ -- заданная функция из $C(G)$ -- \textit{свободный член} (\textit{правая часть}) уравнения.

\bigskip

Здесь $L$ -- линейный дифференциальный оператор второго порядка с постоянными коэффициентами:
\[L=\sum _{i,j=1}^{n}a_{ij} \frac{\partial ^{2} }{\partial x_{i} \partial x_{j} }  +\sum _{i=1}^{n}b_{i} \frac{\partial }{\partial x_{i} }  +c,\] 
преобразующий функции из $C^{2} (G)$ в функции из $C(G)$. 

\bigskip

Запишем оператор $L$ в виде:
\[L=\sum _{i,j=1}^{n}a_{ij} \frac{\partial }{\partial x_{i} } \cdot \frac{\partial }{\partial x_{j} }  +\sum _{i=1}^{n}b_{i} \frac{\partial }{\partial x_{i} }  +c\] 
Первая сумма -- это \textit{квадратичная форма}, вторая -- \textit{линейная форма} от $\frac{\partial }{\partial x_{1} } ,...,\frac{\partial }{\partial x_{n} } $.

\noindent 



Осуществим в ЛДУЧП2ППК  линейную невырожденную замену переменных:
\[{\rm \xi }_{k} =\sum _{i=1}^{n}s_{ik} x_{i}  ,   k=1,...,n\] 
или, иначе, (в матричном виде)
\[{\rm \xi }=\left(\begin{array}{c} {{\rm \xi }_{1} } \\ {...} \\ {{\rm \xi }_{n} } \end{array}\right)=S^{T} \left(\begin{array}{c} {x_{1} } \\ {...} \\ {x_{n} } \end{array}\right)=S^{T} x,   S=(s_{ik} ),   \det S\ne 0\] 
или в развернутом виде
\[\left\{\begin{array}{c} {{\rm \xi }_{1} =s_{11} x_{1} +\ldots +s_{n1} x_{n} } \\ {\ldots } \\ {{\rm \xi }_{k} =s_{1k} x_{1} +\ldots +s_{nk} x_{n} } \\ {\ldots } \\ {{\rm \xi }_{n} =s_{1n} x_{1} +\ldots +s_{nn} x_{n} } \end{array}\right. \] 
Тогда
\[\frac{\partial {\rm \xi }_{k} }{\partial x_{i} } =s_{ik} ,   i,k=1,...,n.\] 

Так как $\det S\ne 0$, то можно выразить переменные $x_{1} ,...,x_{n} $ через переменные ${\rm \xi }_{1} {\rm ,...,\xi }_{n} $: $x=x({\rm \xi })=(S^{-1} )^{T} {\rm \xi }$. Обозначим $u(x({\rm \xi }))=u((S^{-1} )^{T} {\rm \xi })=\tilde{u}({\rm \xi })$, тогда $\tilde{u}(S^{T} x)=u(x)$. Имеем
\[\frac{\partial u}{\partial x_{i} } =\sum _{k=1}^{n}\frac{\partial \tilde{u}}{\partial {\rm \xi }_{k} } \frac{\partial {\rm \xi }_{k} }{\partial x_{i} }  =\sum _{k=1}^{n}\frac{\partial \tilde{u}}{\partial {\rm \xi }_{k} } s_{ik},\quad   \frac{\partial }{\partial x_{i} } =\sum _{k=1}^{n}s_{ik} \frac{\partial }{\partial {\rm \xi }_{k} }  \] 
\[\frac{\partial ^{2} u}{\partial x_{i} \partial x_{j} } =\frac{\partial }{\partial x_{i} } \left(\frac{\partial u}{\partial x_{j} } \right)=\left[\frac{\partial }{\partial x_{i} } \cdot \frac{\partial }{\partial x_{j} } \right]u=\left[\sum _{k=1}^{n}s_{ik} \frac{\partial }{\partial {\rm \xi }_{k} }  \cdot \sum _{l=1}^{n}s_{jl} \frac{\partial }{\partial {\rm \xi }_{l} }  \right]\tilde{u}=\] 
\[=\sum _{k,l=1}^{n}s_{ik} s_{jl} \frac{\partial ^{2} \tilde{u}}{\partial {\rm \xi }_{k} \partial {\rm \xi }_{l} }  \] 
Подставляя найденные выражения для $\frac{\partial u}{\partial x_{i} } $ и $\frac{\partial ^{2} u}{\partial x_{i} \partial x_{j} } $ в  ЛДУЧП2ППК , получим:
\[\sum _{k,l=1}^{n}\frac{\partial ^{2} \tilde{u}}{\partial {\rm \xi }_{k} \partial {\rm \xi }_{l} } \sum _{i,j=1}^{n}a_{ij} s_{ik} s_{jl}   +\sum _{k=1}^{n}\frac{\partial \tilde{u}}{\partial {\rm \xi }_{k} }  \sum _{i=1}^{n}b_{i} s_{ik}  +c\tilde{u}=\tilde{f}({\rm \xi })\] 
или
\[\tilde{L}\tilde{u}=\sum _{k,l=1}^{n}\tilde{a}_{kl} \frac{\partial ^{2} \tilde{u}}{\partial {\rm \xi }_{k} \partial {\rm \xi }_{l} }  +\sum _{k=1}^{n}\tilde{b}_{k} \frac{\partial \tilde{u}}{\partial {\rm \xi }_{k} }  +c\tilde{u}=\tilde{f}({\rm \xi })\] 
где
\[\tilde{a}_{kl} =\sum _{i,j=1}^{n}a_{ij} s_{ik} s_{jl},\quad \tilde{b}_{k} =\sum _{i=1}^{n}b_{i} s_{ik},\quad   \tilde{f}({\rm \xi })=f(x({\rm \xi })),\] 
\[\tilde{L}=\sum _{k,l=1}^{n}\tilde{a}_{kl} \frac{\partial ^{2} }{\partial {\rm \xi }_{k} \partial {\rm \xi }_{l} }  +\sum _{k=1}^{n}\tilde{b}_{k} \frac{\partial }{\partial {\rm \xi }_{k} }  +c\] 



\noindent Запишем выражение для преобразования коэффициентов $a_{ij} $ ($i,j=1,...,n$) при старших (вторых) производных уравнения (ЛДУЧП2ППК) в матричном виде:

\noindent $\tilde{A}=S^{T} AS$, где $A=(a_{ij} )$, $\tilde{A}=(\tilde{a}_{kl} )$.



Полученная формула преобразования старших коэффициентов $a_{ij} $~($i,j=1,...,n$) совпадает с формулой преобразования коэффициентов квадратичной формы

\noindent $K=\sum _{i,j=1}^{n}a_{ij} y_{i} y_{j}  =y^{T} Ay$,   где    $y=\left(\begin{array}{c} {y_{1} } \\ {...} \\ {y_{n} } \end{array}\right)$

\noindent при неособенной линейной замене переменных

\noindent $y_{i} =\sum _{k=1}^{n}s_{ik} {\rm \eta }_{k}  $   (или, иначе, $y=\left(\begin{array}{c} {y_{1} } \\ {...} \\ {y_{n} } \end{array}\right)=S\left(\begin{array}{c} {{\rm \eta }_{1} } \\ {...} \\ {{\rm \eta }_{n} } \end{array}\right)=S{\rm \eta }$),   $\det S\ne 0$,

\noindent приводящей форму $K$ к виду:
\[K=\sum _{k,l=1}^{n}\tilde{a}_{kl} {\rm \eta }_{k} {\rm \eta }_{l}  ={\rm \eta }^{T} \tilde{A}{\rm \eta }\] 

